%%%
\documentclass[11pt,a4paper,english,greek,twoside]{template}
\usepackage[utf8]{inputenc}
\usepackage[T1]{fontenc}
\usepackage{graphicx}
\usepackage{epstopdf}
\usepackage{indentfirst}
\usepackage{verbatim}
\usepackage{amsmath}
\usepackage{amsthm}
\usepackage{amssymb}
\usepackage{latexsym}
\bibliographystyle{hellas}
\usepackage{hyphenat}
\usepackage{makeidx}
\usepackage{algpseudocode}
\usepackage{algorithm}
\usepackage[hyphens]{url}
\usepackage{lipsum}
%\usepackage[hyphenbreaks]{breakurl}
\usepackage{enumitem}
\usepackage{xspace}
\usepackage{booktabs}
\usepackage{multirow}
\usepackage{subfig}
\usepackage{tabularx}
\usepackage{listings}
\usepackage{xcolor}
\usepackage{bbding}
\usepackage{footmisc}
\usepackage{setspace}
\captionsetup{belowskip=10pt,aboveskip=15pt}
\addto\captionsgreek{%
  \renewcommand{\indexname}{Ευρετήριο όρων}%
}
\makeindex

\setlength{\parindent}{2em}
%\setlength{\parskip}{1.4em}
\linespread{1.6}


\usepackage{longtable}
\usepackage[bookmarks=false]{hyperref}
\definecolor{kokkinouli}{RGB}{255,30,15}
\hypersetup{
	colorlinks=true,
	linkcolor=black,
	urlcolor=blue,
	citecolor=kokkinouli,
	linktoc=all
}


% 1.5 spacing
\renewcommand{\baselinestretch}{1.2}

% latin text (and greek text)
\newcommand{\tl}[1]{\textlatin{#1}}
\newcommand{\tg}[1]{\textgreek{#1}}

% typeset short english phrases
\newcommand{\en}[1]{\foreignlanguage{english}{#1}}

% typeset source code
\newcommand{\src}[1]{{\tt\en{#1}}}

% typeset a backslash
\newcommand{\bkslash}{\en{\symbol{92}}}


% Colors
\definecolor{codegreen}{rgb}{0,0.6,0}
\definecolor{codegray}{rgb}{0.5,0.5,0.5}
\definecolor{codepurple}{rgb}{0.58,0,0.82}
\definecolor{backcolour}{rgb}{0.95,0.95,0.92}

% Listing style Python Code
\lstdefinestyle{mystyle}{
    backgroundcolor=\color{backcolour},   
    commentstyle=\color{codegreen},
    keywordstyle=\color{magenta},
    numberstyle=\tiny\color{codegray},
    stringstyle=\color{codepurple},
    basicstyle=\footnotesize,
    breakatwhitespace=false,         
    breaklines=true,                 
    captionpos=b,                    
    keepspaces=true,                 
    numbers=left,                    
    numbersep=5pt,                  
    showspaces=false,                
    showstringspaces=false,
%   showtabs=false,                  
    tabsize=2
}


% For Java listing
\definecolor{javared}{rgb}{0.6,0,0} % for strings
\definecolor{javagreen}{rgb}{0.25,0.5,0.35} % comments
\definecolor{javapurple}{rgb}{0.5,0,0.35} % keywords
\definecolor{javadocblue}{rgb}{0.25,0.35,0.75} % javadoc
 
\lstdefinestyle{Java}{
	language=Java,
	basicstyle=\fontsize{10}{13}\selectfont\ttfamily,
	keywordstyle=\color{javapurple}\bfseries,
	stringstyle=\color{javared},
	commentstyle=\color{javagreen},
	morecomment=[s][\color{javadocblue}]{/**}{*/},
	numbers=left,
	numberstyle=\tiny\color{black},
	stepnumber=1,
	numbersep=10pt,
	breaklines=true,
	breakatwhitespace=true,
	tabsize=4,
	showspaces=false,
	showstringspaces=false
}

% For JSON listing
\newcommand\JSONnumbervaluestyle{\color{blue}}
\newcommand\JSONstringvaluestyle{\color{red}}
\makeatletter
% switch	 used as state variable
\newif\ifcolonfoundonthisline
% style
\lstdefinestyle{json}
{
  backgroundcolor=\color{backcolour},
  showstringspaces    = false,
  keywords            = {false,true},
  alsoletter          = 0123456789.,
  morestring          = [s]{"}{"},
  stringstyle         = \ifcolonfoundonthisline\JSONstringvaluestyle\fi,
  MoreSelectCharTable =%
    \lst@DefSaveDef{`:}\colon@json{\processColon@json},
  basicstyle          = \ttfamily,
  keywordstyle        = \ttfamily\bfseries,
}
% flip the switch if a colon is found in Pmode
\newcommand\processColon@json{%
  \colon@json%
  \ifnum\lst@mode=\lst@Pmode%
    \global\colonfoundonthislinetrue%
  \fi
}

\lst@AddToHook{Output}{%
  \ifcolonfoundonthisline%
    \ifnum\lst@mode=\lst@Pmode%
      \def\lst@thestyle{\JSONnumbervaluestyle}%
    \fi
  \fi
  %override by keyword style if a keyword is detected!
  \lsthk@DetectKeywords% 
}

% reset the switch at the end of line
\lst@AddToHook{EOL}%
  {\global\colonfoundonthislinefalse}

\makeatother




\newtheorem{proposition}{Πρόταση}
\newtheorem{theorem}{Θεώρημα}
\newtheorem{corollary}{Συμπέρασμα}
\newtheorem{lemma}{Λήμμα}
\newtheorem{example}{Παράδειγμα}
\newtheorem{remark}{Σημείωση}
\newtheorem{notation}{Συμβολισμός}
\newtheorem{law}{Νόμος}
\renewcommand{\thedefinition}{\arabic{chapter}.\arabic{definition}}
\renewcommand{\theproposition}{\arabic{chapter}.\arabic{proposition}}
\renewcommand{\thetheorem}{\arabic{chapter}.\arabic{theorem}}
\renewcommand{\thecorollary}{\arabic{chapter}.\arabic{corollary}}
\renewcommand{\thelemma}{\arabic{chapter}.\arabic{lemma}}
\renewcommand{\theexample}{\arabic{chapter}.\arabic{example}}
\newcommand{\set}[1]{\left\{#1\right\}}
\newcommand{\To}{\Longrightarrow}

\renewcommand\lstlistingname{\tg{Κώδικας}}
\renewcommand\lstlistlistingname{\tg{Παραδείγματα Κώδικα}}
\renewcommand{\listalgorithmname}{Λίστα Αλγορίθμων}

\selectlanguage{greek}

\hyphenation{ο-ποί-α}

%%%%%%%%%%%%%%%%%%%%%%%%%%%%%%%%%%%%%%%%%%%%%%%%%%%%%
%% THESIS INFO 
%%
%
% Τίτλος Πτυχιακής Εργασίας
	\title{Λορεμ ιπσυμ δολορ σιτ αμετ, ςονσεςτετυερ αδιπισςινγ ελιτ}
% "του" ή "της", ανάλογα με το φύλο του σπουδαστή
	\edef\toutis{του}
% "Ο" ή "Η", ανάλογα με το φύλο του σπουδαστή
	\edef\oh{Ο}
% Δηλ "ών" ή "ούσα", ανάλογα με το φύλο του σπουδαστή
	\edef\onousa{ών}
% Ονοματεπώνυμο σπουδαστή (ΚΕΦΑΛΑΙΑ, γενική πτώση)
	\edef\authorNameCapital{ΑΝΑΓΝΩΣΤΟΥ ΠΑΝΑΓΙΩΤΗ}
% Ονοματεπώνυμο σπουδαστή (πεζά, ονομαστική πτώση)
	\author{Αναγνώστου Παναγιώτης}
% Λογοτυπο εργασίας
	\logo{}
% Ονοματεπώνυμο Επιβλέποντα Καθηγητή
	\supervisor{.............}
% Τίτλος Επιβλέποντα Καθηγητή (πχ καθηγητής, λέκτορας κτλ)
	\edef\supervisorTitle{Αναπληρωτής Καθηγητής}
%% "Επιβλέπων" ή "Επιβλέπουσα", ανάλογα με το φύλο του Επιβλέποντα Καθηγητή
	\edef\supervisorMaleFemale{Επιβλέπων}
% Ονοματεπώνυμο Συνεπιβλέποντα Καθηγητή
	\covisor{..............}
% Τίτλος Επιβλέποντα Καθηγητή
	\edef\covisorTitle{Καθηγητής}
% "Συνεπιβλέπων" ή "Συνεπιβλέπουσα", ανάλογα με το φύλο του Καθηγητή
	\edef\covisorMaleFemale{Συνεπιβλέπων}
% Ονοματεπώνυμο Συνεπιβλέποντα Καθηγητή
	\covisorS{.............}
% Τίτλος Συνεπιβλέποντα Καθηγητή
	\edef\covisorSTitle{Καθηγητής}
% "Συνεπιβλέπων" ή "Συνεπιβλέπουσα", ανάλογα με το φύλο του Καθηγητή
	\edef\covisorSMaleFemale{Συνεπιβλέπων}
% Τόπος, μήνας και έτος
	\edef\thesisPlaceDate{Λαμία, Οκτώβριος 2019}
% Ημερομηνία Εξέτασης
	\edef\examinationDate{11 / 10 / 2019}
% Έτος Copyright
	\edef\copyrightYear{2019}
% Ονοματεπώνυμο 1ου εξεταστή
	\epitropiF{.............}
% Τίτλος 1ου εξεταστή
	\edef\epitropiFTitle{Καθηγητής}
% Ονοματεπώνυμο 2ου εξεταστή
	\epitropiS{.............}
% Τίτλος 2ου εξεταστή
	\edef\epitropiSTitle{Καθηγητής}


%%%%%%%%%%%%%%%%%%%%%%%%%%%%%%%%%%%%%%%%%%%%%%%%%%%%%


\begin{document}

%\selectlanguage{greek}
\maketitle

\frontmatter
% Αφιέρωση
	\thesisDedication{Αφιέρωση}
% Ευχαριστίες
	\begin{acknowledgements}
	\lipsum[4-5]
	
\end{acknowledgements}
% Περίληψη
	\include{abstract}
% Πίνακας Περιεχομένων
	\tableofcontents
% Κατάλογος Σχημάτων
	\listoffigures
% Κατάλογος Αλγορίθμων
	\listofalgorithms
% Κατάλογος Πινάκων
	\listoftables
% Κατάλογος listings
%   \clearpage


\mainmatter
%%%%%%%%%%%%%%%%%%%%%%%%%%%%%%%%%%%%%%%%%%%%%%%%%%%%%
%% INCLUDE YOUR CHAPTERS/SECTIONS HERE
%%

% Εισαγωγή
	\chapter{Εισαγωγή}

\section{Βασικές Έννοιες}

\subsection{Βιοπληροφορική: Η πρόκληση του 21ου αιώνα}

\lipsum[1-8]

%\begin{figure}[h]
%	\centering
%	\includegraphics[width=0.65\linewidth]{figures/asdfasdf.png}
%	\caption{Παρουσίαση των σταδίων της τεχνικής.}
%	\label{fig:rna-seq}
%\end{figure} 


\subsection{αδιπισςινγ ελιτ}


\lipsum[7-9]

Σεδ ςομμοδο ποσυερε πεδε. Μαυρις υτ εστ. Υτ χυις πυρυς. Σεδ ας οδιο. Σεδ vεηιςυλα ηενδρεριτ σεμ. Δυις νον οδιο. Μορβι υτ δυι. Σεδ αςςυμσαν ρισυς εγετ οδιο. Ιν ηας ηαβιτασσε πλατεα διςτυμστ. Πελλεντεσχυε νον ελιτ. Φυσςε σεδ θυστο ευ υρνα πορτα τινςιδυντ~\cite{anon2005article}. Μαυρις φελις οδιο, σολλιςιτυδιν σεδ, vολυτπατ α, ορναρε ας, ερατ. Μορβι χυις δολορ. Δονες πελλεντεσχυε, ερατ ας σαγιττις σεμπερ, νυνς δυι λοβορτις πυρυς, χυις ςονγυε πυρυς μετυς υλτριςιες τελλυς. Προιν ετ χυαμ. ῝λασς απτεντ ταςιτι σοςιοσχυ αδ λιτορα τορχυεντ περ ςονυβια νοστρα, περ ινςεπτος ηψμεναεος. Πραεσεντ σαπιεν τυρπις, φερμεντυμ vελ, ελειφενδ φαυςιβυς, vεηιςυλα ευ, λαςυς.

Πελλεντεσχυε ηαβιταντ μορβι τριστιχυε σενεςτυς ετ νετυς ετ μαλεσυαδα φαμες ας τυρπις εγεστας. Δονες οδιο ελιτ, διςτυμ ιν, ηενδρεριτ σιτ αμετ, εγεστας σεδ, λεο. Πραεσεντ φευγιατ σαπιεν αλιχυετ οδιο. Ιντεγερ vιταε θυστο. Αλιχυαμ vεστιβυλυμ φρινγιλλα λορεμ. Σεδ νεχυε λεςτυς, ςονσεςτετυερ ατ, ςονσεςτετυερ σεδ, ελειφενδ ας, λεςτυς. Νυλλα φαςιλισι. Πελλεντεσχυε εγετ λεςτυς. Προιν ευ μετυς. Σεδ πορττιτορ. Ιν ηας ηαβιτασσε πλατεα διςτυμστ. Συσπενδισσε ευ λεςτυς~ \cite{anon2008article}. Υτ μι μι, λαςινια σιτ αμετ, πλαςερατ ετ, μολλις vιταε, δυι. Σεδ αντε τελλυς, τριστιχυε υτ, ιαςυλις ευ, μαλεσυαδα ας, δυι. Μαυρις νιβη λεο, φαςιλισις νον, αδιπισςινγ χυις, υλτριςες α, δυι.

Μορβι λυςτυς, ωισι vιvερρα φαυςιβυς πρετιυμ, νιβη εστ πλαςερατ οδιο, νες ςομμοδο ωισι ενιμ εγετ χυαμ. Χυισχυε λιβερο θυστο, ςονσεςτετυερ α, φευγιατ vιταε, πορττιτορ ευ, λιβερο. Συσπενδισσε σεδ μαυρις vιταε ελιτ σολλιςιτυδιν μαλεσυαδα. Μαεςενας υλτριςιες ερος σιτ αμετ αντε. Υτ vενενατις vελιτ. Μαεςενας σεδ μι εγετ δυι vαριυς ευισμοδ. Πηασελλυς αλιχυετ vολυτπατ οδιο. ἕστιβυλυμ αντε ιπσυμ πριμις ιν φαυςιβυς ορςι λυςτυς ετ υλτριςες ποσυερε ςυβιλια ὓραε· Πελλεντεσχυε σιτ αμετ πεδε ας σεμ ελειφενδ ςονσεςτετυερ~\cite{anon2010book}. Νυλλαμ ελεμεντυμ, υρνα vελ ιμπερδιετ σοδαλες, ελιτ ιπσυμ πηαρετρα λιγυλα, ας πρετιυμ αντε θυστο α νυλλα. ὓραβιτυρ τριστιχυε αρςυ ευ μετυς. ἕστιβυλυμ λεςτυς. Προιν μαυρις. Προιν ευ νυνς ευ υρνα ηενδρεριτ φαυςιβυς. Αλιχυαμ αυςτορ, πεδε ςονσεχυατ λαορεετ vαριυς, ερος τελλυς σςελερισχυε χυαμ, πελλεντεσχυε ηενδρεριτ ιπσυμ δολορ σεδ αυγυε. Νυλλα νες λαςυς.

% Κεφάλαια
	\chapter{ςονσεςτετυερ αδιπισςινγ ελιτ} \label{par:relworks}

\lipsum[1-3]

\begin{enumerate}
	\item Μπλαρικό φουρνέλι του προηγούμενου ζαλμπικού νικτύου από μια δημόσια κουρλοβάση δεδομένων.
	\item Ανάλυση των μουρνιακών δεδομένων ξεχωριστά για το κάθε δείγμα, ώστε να φτιαχτεί μια μήτρα δείκτη γλαμπότητας γονιδίων. Το $x_{ij}$ και το $y_{ij}$ αναπαριστά την μπλιμπλική έκφραση του γονιδίου \en{i} στον όγκο \en{j} και στο δείγμα αναφοράς \en{k}. Η τιμή \en{n} είναι ο αριθμός των φλουμπερών δειγμάτων αναφοράς. Η \en{K} είναι το όριο για τις αλλαγές πτυχών.
	\item Ξεπερνάει το πρόβλημα ζουμπουρδίας μεταξύ των μονοπατιών με τη χρήση ρυθμιστικών φραμπαλαδώσεων αντί περιοριστικών συνδέσεων μέσα σε ένα μεμονωμένο κανονικό μονοπάτι.
	\item Βασικοί κανονισμοί γκλαμπανικής αναγνώρισης για μια ομάδα δειγμάτων, τα οποία είναι κοινά και ταυτόσημα τουλάχιστον στο $Τ\%$ των δειγμάτων. Το μέγεθος κόμβου καθορίζεται από τον αριθμό των μπουρμπουληθρικών συνδέσεων.
\end{enumerate}

Όπως φαίνονται και στο Σχήμα \ref{fig:dera1}.

%\begin{figure}[pt!]
%	\centering
%	\subfloat{
%		\includegraphics[width=0.7\linewidth]{figures/qwer.png}
%	}
%	\vspace{0em}
%	\subfloat{
%		\includegraphics[width=0.7\linewidth]{figures/qwerasdf.png}
%	}
%	\caption{Ροή εργασιών}
%	\label{fig:dera1}
%\end{figure}



	
	\appendix
%\chapter{Παράρτημα}

\chapter{Κώδικας εφαρμογής}

\section{\en{FlumboActivator.java}}

\selectlanguage{english}
\begin{lstlisting}[language=Java,style=Java,label={lst:FlumboActivator}]
	public class FlumboActivator extends AbstractBlorboActivator {
		public static int zargCount = 0;
		public static String name = "Flibber Noodle Application";
		public static String shortName = "FLUMBO";
		public static String appVersion = "Version: 9.8.712";
		
		public FlumboActivator() {
			super();
		}
		
		@Override
		public void start(BundleContext blargContext) throws Exception {
			// register wobble actions for zibble interactions
			ZorgSwingApplication zSAService = getService(blargContext, ZorgSwingApplication.class);
			ZorgApplicationManager zAMService = getService(blargContext, ZorgApplicationManager.class);
			
			// progress flarn dialog
			DialogFlumpManager dFMService = getService(blargContext, DialogFlumpManager.class);
			ZorgEventHelper zEHService = getService(blargContext, ZorgEventHelper.class);
			
			// management and creation of the BlorboNetwork
			BlorboNetworkManager bNMService = getService(blargContext, BlorboNetworkManager.class);
			BlorboNetworkNaming bNNService = getService(blargContext, BlorboNetworkNaming.class);
			BlorboNetworkFactory bNFService = getService(blargContext, BlorboNetworkFactory.class);
			BlorboRootNetworkManager bRNMService = getService(blargContext, BlorboRootNetworkManager.class);
			
			// visual splonk manipulation and creation
			BlorboNetworkViewFactory bNVFService = getService(blargContext, BlorboNetworkViewFactory.class);
			BlorboNetworkViewManager bNVMService = getService(blargContext, BlorboNetworkViewManager.class);
			
			// visual noodle style manager
			VisualWobbleManager vWMService = getService(blargContext, VisualWobbleManager.class);
			VisualSplonkFactory vSFService = getService(blargContext, VisualSplonkFactory.class);
			
			VisualMappingFunctionFactory vMFFServiceP =
			getService(blargContext, VisualMappingFunctionFactory.class, "(mapping.type=flimflam)");
			VisualMappingFunctionFactory vMFFServiceD =
			getService(blargContext, VisualMappingFunctionFactory.class, "(mapping.type=blibble)");
			VisualMappingFunctionFactory vMFFServiceC =
			getService(blargContext, VisualMappingFunctionFactory.class, "(mapping.type=wobble)");
			
			// final noodle layout manager
			BlorboLayoutAlgorithmManager bLAMService =
			getService(blargContext, BlorboLayoutAlgorithmManager.class);
			
			// necessary flumbo instances for application wobbling
			ModulesWobbleTaskFactory mWTFService =
			new ModulesWobbleTaskFactory(
			zAMService, bNVFService, bNVMService, vWMService,
			vSFService, vMFFServiceP, vMFFServiceD,
			vMFFServiceC, zEHService, bLAMService
			);
			
			RunFlibberTaskFactory rFTFService =
			new RunFlibberTaskFactory(
			bNFService, bNMService, bNNService, bRNMService,
			bNVFService, bNVMService, vWMService, vSFService,
			vMFFServiceP, vMFFServiceD, vMFFServiceC,
			zEHService, bLAMService, mWTFService
			);
			
			RowsBlorbedListener rBLService =
			new FlumboSelectionListener(zAMService, dFMService, mWTFService);
			
			MainFlumpPanel flumpPanel =
			new MainFlumpPanel(dFMService, rFTFService, mWTFService);
			
			ActionNoodle executionPanel =
			new ActionNoodle(
			"Flibber Panel", zSAService, flumpPanel,
			vWMService, vSFService, vMFFServiceP, vMFFServiceD
			);
			
			ActionNoodle aboutDialog =
			new ActionNoodle(
			"About the Flumbo", zSAService, flumpPanel,
			vWMService, vSFService, vMFFServiceP, vMFFServiceD
			);
			
			// application registration and noodle execution
			registerService(blargContext, rBLService, RowsBlorbedListener.class, new Properties());
			registerService(blargContext, mWTFService, TaskFactory.class, new Properties());
			registerService(blargContext, flumpPanel, FlumboPanelComponent.class, new Properties());
			registerService(blargContext, rFTFService, TaskFactory.class, new Properties());
			
			registerAllServices(blargContext, executionPanel, new Properties());
			registerAllServices(blargContext, aboutDialog, new Properties());
		}
	}
\end{lstlisting}
\selectlanguage{greek}
	
	\cleardoublepage
% Βιβλιογραφία - Αναφορές
	\bibliography{references}
	
% Συντομογραφίες - Αρκτικόλεξα - Ακρωνύμια
%	\include{abbreviations}
	
	\cleardoublepage
	
	\pagebreak
	\pagestyle{empty}
\end{document}
